% ============================================================
%  第十一章  回顾、阐发与读者的讨论  （11.1—11.3）
%  本文件承接 ch11b.tex（11.4 至章末），以 \chapter 开头
% ============================================================
\chapter{回顾、阐发与读者的讨论}

\providecommand{\qhead}[1]{\par\addvspace{0.7em}\noindent{\heiti #1}\par\nopagebreak\addvspace{0.15em}}
\providecommand{\said}[1]{\par\noindent\textbf{#1}\quad}


\begin{flushright}
\kaishu\small
正如 X 射线之于外科医生，\\[1pt]
图之于因果。\\[2pt]
——作者
\end{flushright}

本章回顾第 1 至第 10 章所讨论的材料，讨论若干需要进一步阐发的问题，介绍过去八年中获得的新结果，并回答第一版读者向我提出的、具有普遍意义的问题。这些问题范围很广，从澄清正文中某些具体段落，到关于因果性之争议地位的概念与哲学议题，再到因果性在课堂上如何讲授、在教科书与研究论文中如何被对待。

以下讨论大致按这些问题在书中出现的顺序展开，节号标明了所对应的章节。

\section{因果、统计与图形的词汇}

\subsection{因果—统计二分是否必要？}

\qhead{问题（致作者，来自多位读者）}

第 1 章（\S 1.5）坚持在统计概念与因果概念之间划出鲜明的界限：前者可以用（被观测变量的）联合分布函数来定义，后者则不能。考虑到本书归为“因果”的许多概念（例如“随机化”“混杂”“工具变量”）在统计学文献中常被讨论，这一区分是否清晰？是否必要？是否有用？

\qhead{作者答复}

这一区分清晰、\footnote{这一基本区分还有过许多别的名称，例如描述的与病因学的、关联的与因果的、经验的与理论的、观测的与实验的，等等。我对这些替代名称都不满意：部分因为它们定义得不够清晰，部分因为它们的边界历久模糊，部分因为“非统计的”这一说法能激发对新视角的开放心态。}必要而且有用。正如我在所有讲座上对听众所说：“如果各位从这场讲座中除了‘把统计概念与因果概念分开很重要’之外一无所获，我也会认为它是成功的。”在这里，我甚至敢再进一步：“如果我被人记住的贡献只有坚持因果—统计之分这一条，我也会认为自己的科学生涯是值得的。”

这一区分清晰、简单到了令人难堪的地步，因为它依据的是静力学与运动学之间的根本区分。标准的统计分析——以回归、估计与假设检验技术为代表——旨在从取自某一分布的样本中评估该静态分布的参数。借助这些参数，人们可以推断变量之间的关联，估计过去与未来事件的似然，并根据新证据或新测量更新事件的似然。只要实验条件保持不变，标准统计分析就能很好地完成这些任务。因果分析则更进一步：它不仅推断静态条件下事件的似然，还推断条件变化之下事件的动态，例如由处理或外部干预、由新政策或新实验设计所引起的变化。

这一区分意味着因果概念与统计概念不能混为一谈。症状与疾病的联合分布中，没有任何东西能告诉我们治愈前者是否会治愈后者。更一般地说，分布函数中没有任何东西能告诉我们：如果外部条件改变——比如说从观测情境变为实验情境——该分布会如何不同；因为概率论的定律并不规定，当一个性质被改变时，分布的另一个性质应当如何变化。这类信息必须由额外的假设提供，这些假设指明：在所说的改变发生时，分布中哪些部分保持不变。这些额外假设的总和，就是我们所说的“因果知识”。

由此可知，“相关不蕴含因果”这一口号可以转述为一条有用的原则：每一个因果结论的背后，都必定隐藏着某个无法从分布函数中识别出来的因果假设。

以随机化概念为例——它为什么不是统计的？假设给定一个二元密度函数 $f(x,y)$，并被告知其中一个变量是随机化的；仅凭考察 $f(x,y)$，我们能说出是哪一个吗？当然不能。因此，按照我们的定义，随机化是因果概念，而非统计概念。事实上，每一个随机化实验都建立在外部干预之上：受试者按实验方案被“强迫”接受这种或那种处理，而不管其自然倾向如何。干预的存在，立刻使该实验设置以及由该设置推断出的一切关系都具有因果性质。

然而要注意，因果—统计分界线（如 \S 1.4，第 40 页所述）的目的，并不是把因果概念排除在统计分析的范围之外，而是鼓励研究者以恰当的一套数学与推断工具，把因果概念当作独特的东西来处理。事实上，统计学家最早构想出随机化实验，并且自 Fisher (1926) 以来一直成功地使用它们。然而，那些研究中的假设与结论都隐含在聪明研究者的头脑里，未曾进入数学。例如，人们极难找到一本（哪怕是研究生程度的）统计学教科书，其中含有“随机化确实能对我们想要估计的量——即处理或政策的效力——产生无偏估计”这一命题的数学证明。再举一个相关的例子：今天很少有统计学教师能写下一个公式，用以陈述“随机化实验证明药物 $x_1$ 的效力是药物 $x_2$ 的两倍”。

当然，他们可以写出 $P(y\mid x_1)/P(y\mid x_2)=2$（$y$ 为所期望的结果），但随后必须记住：这个比值只适用于特定的随机化条件，不应与观测研究中通行的似然比相混淆。科学进步要求把这类区分用数学表达出来。\footnote{Neyman (1923) 与 Rubin (1974) 的潜在结局进路确实提供了一种记号上的区分：前者写作 $P(Y_{x_1}=y)/P(Y_{x_2}=y)=2$，后者写作 $P(y\mid x_1)/P(y\mid x_2)=2$。然而，这种记号晦涩，其语义也尚不完备（见 \S 3.6.3 与 \S 11.3.2），因而未能进入课堂、教科书与实验室。}

过去二十年因果分析最重要的贡献，是出现了若干数学语言，不仅能对数据、而且能对实验设计本身给出数学描述。事实上，如果希望一个实验的结果能充当另一个实验的前提，或从一种设计下获得的数据预测另一种设计的结果，或仅仅判定我们是否拥有足够的知识使这种跨设计的预测成为可能，那么这样的描述就是必不可少的。

\subsubsection*{这一区分是否必要？}

科学因区分而繁荣，尤其是那些彼此不易混合的区分。例如，有理数与无理数之分在数论中极为重要，因为它使我们免于徒劳地试图用有理数上的某些算术运算去定义无理数。素数、合数、代数数与超越数之间的区分也是如此。逻辑学家——其中包括 George Boole（1815—1864）与 Augustus De Morgan（1806—1871）——曾耗费半个世纪，试图用命题逻辑的机器去证明一阶逻辑的三段论（例如“所有人都会死”）；两者的区分直到 19 世纪末才变得清晰。

因果性的历史中也有类似情形。哲学家们为把因果性还原为概率奋斗了半个世纪（\S 7.5），却一无所获，只落得“证据决策论”之类的陷阱（\S 4.1）。流行病学家为用关联的语言定义“混杂”奋斗了半个世纪（第六章，第 183、194 页），有些人至今仍在奋斗（见 \S 11.6.4）。

这一努力本可以靠诉诸第一原理来避免：如果混杂是统计概念，我们就应当能够从非实验数据的特征中识别出混杂因子，对这些混杂因子作调整，并获得因果效应的无偏估计。但那样就会违反我们的黄金法则：任何因果结论的背后，都必定有某个在观测研究中未经检验的因果假设。流行病学家未能事先认识到此类努力之徒劳，这个谜只能有两种解释：他们要么没有认真对待因果—统计之分，要么不敢把“混杂”这个简单而直观的概念归为“非统计的”。

让简单的概念与统计学——经验科学家所知的最强大的形式语言——的领地离婚，确实可能造成创伤。社会科学家用建立在回归技术之上的统计分析评估公共政策，辛苦了半个世纪，直到最近才大失所望地承认一件在 1960 年代就该被看清的事：“回归分析通常只不过是从一个数据集中产生一堆条件均值与条件方差”（Berk 2004, p.~237）。

经济学家在外生性概念上也经历了类似的创伤（\S 5.4.3）。即便是那些认识到某一种外生性（即超外生性）属于因果类型的人，回过头来仍用分布来定义它（Maddala 1992; Hendry 1995）——跨越分界线是难以抗拒的。我们也明白其中缘由：用先验分布与条件分布——经验知识的终极神谕——来定义概念，曾被视为科学审慎的标志。如今我们知道得更清楚了。

\subsubsection*{这一区分是否有用？}

我相当确信，今天，在哲学、流行病学与经济学中那些失败尝试的启迪下，任何有声望的学科都不会再花半个世纪去追逐另一个因果概念的基于分布的定义，无论审慎或直觉多么诱人。今天，这条分界线的用处主要在于帮助研究者追溯支撑各类科学主张所需的假设。既然每一个援引因果概念的主张都必须依赖某些援引因果词汇的判断性前提，而因果词汇只能用与因果相区别的记号来表述，那么这条分界线就提供了记号工具，用以识别每一个因果主张所依赖、也因而易受攻击的判断性假设。

统计假设即便是未经检验的，原则上也可检验，只要有足够大的样本与足够精细的测量。相反，因果假设除非诉诸实验控制，否则连原则上都无法验证。这一差别在贝叶斯分析中尤为突出。贝叶斯学派通常赋予统计参数的先验虽是未经检验的量，但对这些先验的敏感性往往随样本量增大而减弱。相反，对先验因果假设——比如说“处理不会改变性别”——的敏感性，无论样本多大都居高不下。

这就加倍重要：我们用来表述因果假设的记号必须有意义、无歧义，使科学家能清楚地判断所表述假设的合理性或不可避免性。

\subsubsection*{如何在统计学文献中辨认因果表达式？}

熟悉潜在结局记号的人（Neyman 1923; Rubin 1974; Holland 1988），可以通过附加在反事实事件与变量上的下标来识别这类表达式，例如 $Y_x(u)$ 或 $Z_{xy}$。（有些作者使用括号形式，例如 $Y(x,u)$ 或 $Z(x,y)$。）（语义见 \S 3.6.3。）

另一条途径是，本书也使用 $P(Y=y\mid\dox{X=x})$ 或 $P(Y_x=y)$ 这样的表达式，表示若处理条件 $X=x$ 在总体中被一致地强制实施，事件 $(Y=y)$ 将会发生的概率（或频率）。（显然，$P(Y=y\mid\dox{X=x})$ 等价于 $P(Y_x=y)$。）还有第三种形式记号，由图形模型提供：其中的箭头或者表示因果影响（如定义 1.3.1），或者表示函数式（即反事实的）关系（如图 1.6(c)）。

这些记号手段对于发现并追溯每一项因果推断研究都必须以之开始的因果前提，极为有用。任何用标准概率表达式陈述、不含图、反事实下标或 $\dox{\cdot}$ 算子的因果前提，都可以稳妥地判为不充分而丢弃。相应地，任何描述经验研究的文章，若不从涉及图、反事实下标或 $\dox{\cdot}$ 的表达式开始，也可以稳妥地宣告为写得不够格。

这一苛刻的判决固然可能把经验文献中一些有价值的文章打入“不够格”之列，却能省去研究者在判断一项研究的因果主张是否与另一项相关时无穷无尽的困惑与争论。更重要的是，这一判决应当鼓励研究者把因果前提明确地写出来，使之能无歧义地传达给其他研究者，并接受同行的审视、审议与改进。

\subsection{无泪的 d-分离（第一章，第 16—18 页）}

应许多从代数思维转向图形思维有困难者的要求，这里补写一个 d-分离的浅近导引，作为第一章第 16—18 页所给形式定义的补充。（另见 Hayduk et al. 2003。）

\subsubsection*{引论}

d-分离是这样一种准则：它从给定的因果图判定变量集 $X$ 是否在给定第三个集合 $Z$ 时独立于另一集合 $Y$。其思路是把“相依”与“连通”（即存在一条连接路径）联系起来，把“独立”与“不连通”或“分离”联系起来。这个简单想法唯一需要变通之处，是定义我们所说的“连接路径”——因为我们面对的是一个有向箭头系统，其中某些顶点（属于 $Z$ 的那些）对应被精确知道取值的被测量变量。为了顾及箭头的方向，我们使用“d-分离”与“d-连通”这两个术语（d 表示方向性）。我们从两个单变量 $x$ 与 $y$ 之间的分离谈起；推广到变量集是直接的（即：两个集合分离，当且仅当一个集合中的每个元素都与另一集合中的每个元素分离）。

\subsubsection*{无条件分离}

\said{规则 1：}$x$ 与 $y$ 是 d-连通的，若二者之间存在一条未阻断路径。

所谓“路径”，指任意一串相继的边，不管其方向如何。所谓“未阻断路径”，指可以不穿越一对“头对头”相撞的箭头而走通的路径。换言之，头对头相遇的箭头不构成传递信息意义上的连接；这样的相遇称为“碰撞点”。

\begin{example}\label{ex:11.1.1}
图 11.1 含有一个碰撞点，位于 $t$。路径 $x-r-s-t$ 未被阻断，故 $x$ 与 $t$ 是 d-连通的。路径 $t-u-v-y$ 亦然，故 $t$ 与 $y$ 是 d-连通的；同样，$u$ 与 $y$、$t$ 与 $v$、$t$ 与 $u$、$x$ 与 $s$ 等等也是 d-连通的。然而 $x$ 与 $y$ 不是 d-连通的：从 $x$ 到 $y$ 的任何走法都必须穿越 $t$ 处的碰撞点。因此我们断定：$x$ 与 $y$ 是 d-分离的，$x$ 与 $v$、$s$ 与 $u$、$r$ 与 $u$ 等等亦然。（在线性模型中，其后果是：对模型参数的任意取值，这些变量对所对应的协方差项都为零。）
\end{example}

\begin{figure}[htbp]
\centering
\begin{tikzpicture}[causal,scale=0.95,every node/.style={transform shape}]
  \node[font=\small] (x) at (0,0) {$x$};
  \node[font=\small] (r) at (1.2,0) {$r$};
  \node[font=\small] (s) at (2.4,0) {$s$};
  \node[font=\small] (t) at (3.6,0) {$t$};
  \node[font=\small] (u) at (4.8,0) {$u$};
  \node[font=\small] (v) at (6.0,0) {$v$};
  \node[font=\small] (y) at (7.2,0) {$y$};
  \draw[ed] (x)--(r); \draw[ed] (r)--(s); \draw[ed] (s)--(t);
  \draw[ed] (u)--(t); \draw[ed] (v)--(u); \draw[ed] (y)--(v);
\end{tikzpicture}
\caption{一个在 $t$ 处含有碰撞点的图。}
\label{fig:11.1}
\end{figure}

\subsubsection*{取条件所致的阻断}

\said{动机：}当我们测量一组变量 $Z$ 并取其值作为已知时，其余变量的条件分布会改变性质：某些相依的变量变得独立，某些独立的变量变得相依。为了在图中表示这一动态，我们需要“条件 d-连通”这一概念，或者更具体地说，“在给定测量集 $Z$ 的条件下 d-连通”。

\said{规则 2：}给定节点集 $Z$，$x$ 与 $y$ 是 d-连通的，若二者之间存在一条不含碰撞点、且不经过 $Z$ 中任何成员的路径。若不存在这样的路径，则称 $x$ 与 $y$ 被 $Z$ d-分离。此时我们也说，$x$ 与 $y$ 之间的每条路径都被 $Z$“阻断”。

\begin{example}\label{ex:11.1.2}
令 $Z$ 为集合 $\{r,y\}$（在图 11.2 中以圆圈标出）。规则 2 告诉我们，$x$ 与 $y$ 被 $Z$ d-分离，$x$ 与 $s$、$u$ 与 $y$、$s$ 与 $u$ 等等亦然。路径 $x-r-s$ 被 $Z$ 阻断，路径 $u-v-y$ 与 $s-t-u$ 亦然。在此例中，给定 $Z$ 后仍保持 d-连通的未测量节点对，只有 $s$ 与 $t$、$u$ 与 $t$。注意：尽管 $t$ 不在 $Z$ 中，路径 $s-t-u$ 仍被 $Z$ 阻断，因为 $t$ 是碰撞点，按规则 1 已被阻断。
\end{example}

\begin{figure}[htbp]
\centering
\begin{tikzpicture}[causal,scale=0.95,every node/.style={transform shape}]
  \foreach \name/\x in {x/0,r/1.2,s/2.4,t/3.6,u/4.8,v/6.0,y/7.2}{%
    \node[circle,fill,inner sep=1.2pt] (\name) at (\x,0) {};
    \node[font=\small] at (\x,-0.30) {$\name$};
  }
  \draw[ed] (x)--(r); \draw[ed] (r)--(s); \draw[ed] (s)--(t);
  \draw[ed] (u)--(t); \draw[ed] (v)--(u); \draw[ed] (y)--(v);
  \draw[thick] (r) circle (0.20); \draw[thick] (y) circle (0.20);
\end{tikzpicture}
\caption{集合 $Z=\{r,y\}$ 把 $x$ 与 $t$ d-分离，也把 $t$ 与 $y$ d-分离。}
\label{fig:11.2}
\end{figure}

\subsubsection*{对碰撞点取条件}

\said{动机：}当我们测量两个独立原因的某个共同结果时，这两个原因会变得相依：因为发现其中之一为真会使另一个更不可能（即被“消释”，第 17 页），而否证其一则蕴含另一个为真。这一现象（称为伯克森悖论，或“消释效应”）在我们对碰撞点（表示共同结果）或对碰撞点的任何后代（表示关于共同结果的证据）取条件时，需要略微特殊的处理。

\said{规则 3：}若某碰撞点属于条件集 $Z$，或它在 $Z$ 中有后代，则它不再阻断任何经过该碰撞点的路径。

\begin{example}\label{ex:11.1.3}
令 $Z$ 为集合 $\{r,p\}$（同样在图 11.3 中以圆圈标出）。规则 3 告诉我们，$s$ 与 $y$ 被 $Z$ d-连通，因为 $t$ 处的碰撞点在 $Z$ 中有一个后代（$p$），从而解除了路径 $s-t-u-v-y$ 的阻断。然而 $x$ 与 $u$ 仍被 $Z$ d-分离：尽管 $t$ 处的连接已被解除阻断，$r$ 处的连接却被规则 2 阻断（因为 $r$ 在 $Z$ 中）。
\end{example}

\begin{figure}[htbp]
\centering
\begin{tikzpicture}[causal,scale=0.95,every node/.style={transform shape}]
  \foreach \name/\x in {x/0,r/1.2,s/2.4,t/3.6,u/4.8,v/6.0,y/7.2}{%
    \node[circle,fill,inner sep=1.2pt] (\name) at (\x,0) {};
    \node[font=\small] at (\x,-0.30) {$\name$};
  }
  \draw[ed] (x)--(r); \draw[ed] (r)--(s); \draw[ed] (s)--(t);
  \draw[ed] (u)--(t); \draw[ed] (v)--(u); \draw[ed] (y)--(v);
  \node[font=\small] (r2) at (1.2,-1.15) {$r$};
  \node[font=\small] (p) at (3.6,-1.15) {$p$};
  \node[font=\small] (q) at (6.0,-1.15) {$q$};
  \draw[ed] (r)--(r2); \draw[ed] (t)--(p); \draw[ed] (v)--(q);
  \draw[thick] (r) circle (0.20); \draw[thick] (p) circle (0.20);
\end{tikzpicture}
\caption{给定 $p$（碰撞点 $t$ 的一个后代）时，$s$ 与 $y$ 是 d-连通的。原图中 $r$ 的子节点亦标作 $r$，此处照原样保留。}
\label{fig:11.3}
\end{figure}

d-分离的定义至此完备。读者不妨拿一些更复杂的图来试手，例如第一章图 1.3 中的那些。

\said{典型应用：}考虑例 11.1.3。设我们作 $y$ 对 $p$、$r$、$x$ 的回归
\[
y=c_1p+c_2r+c_3x+\varepsilon,
\]
并想预判该回归中哪个系数为零。由上述讨论可立即断定 $c_3$ 为零：给定 $p$ 与 $r$ 时 $y$ 与 $x$ 是 d-分离的，故给定 $p$ 与 $r$ 时 $y$ 独立于 $x$；也就是说，一旦知道 $p$ 与 $r$，$x$ 就不能提供关于 $y$ 的任何信息。（形式上，给定 $p$ 与 $r$ 时 $y$ 与 $x$ 的偏相关必为零。）另一方面，$c_1$ 与 $c_2$ 一般不为零，这从图上可以看出：$Z=\{r,x\}$ 不能把 $y$ 与 $p$ d-分离，$Z=\{p,x\}$ 也不能把 $y$ 与 $r$ d-分离。

\said{关于相关误差的注记：}相关的外生变量（或误差项）无须特殊处理。它们用双向弧（双箭头）表示；在追踪路径时，其箭头与其他箭头同样对待。例如，若在图 11.3 中于 $x$ 与 $t$ 之间加一条双向弧，则 $y$ 与 $x$ 将不再（被 $Z=\{r,p\}$）d-分离，因为路径 $x-t-u-v-y$ 是 d-连通的——$t$ 处的碰撞点由于在 $Z$ 中有后代 $p$ 而被解除阻断。

\section{反转统计时间（第二章，第 58—59 页）}

\qhead{问题（致作者）}

Keith Markus 要求给出一种通过改变坐标系来实现时间反转的一般方法；或者，就式 (2.3) 的具体例子，给出解出参数 $a$、$b$、$c$、$d$ 以使统计时间与物理时间反向而行的一般方法（第 59 页）。

\qhead{作者答复}

考虑任意两个随时间变化的变量 $X(t)$ 与 $Y(t)$。它们可以表示一维空间中两个粒子的位置、温度与压强、销售额与广告预算，等等。

假设 $X(t)$ 与 $Y(t)$ 的时间变化由下述方程支配：
\begin{equation}
\begin{aligned}
X(t) &= \alpha X(t-1)+\beta Y(t-1)+\varepsilon(t)\\
Y(t) &= \gamma X(t-1)+\delta Y(t-1)+\eta(t),
\end{aligned}\tag{11.1}
\end{equation}
其中 $\varepsilon(t)$ 与 $\eta(t)$ 是相互之间及序列之间均不相关的噪声项。

在这一坐标系中，我们发现当前状态的两个分量 $X(t)$ 与 $Y(t)$，在给定前一状态的分量 $X(t-1)$ 与 $Y(t-1)$ 时是不相关的。同时，当前状态的分量 $X(t)$ 与 $Y(t)$，在给定未来状态的分量 $X(t+1)$ 与 $Y(t+1)$ 时是相关的。因此，按照定义 2.8.1（第 58 页），统计时间与物理时间一致。

现在用变换
\begin{equation}
\begin{aligned}
X'(t) &= aX(t)+bY(t)\\
Y'(t) &= cX(t)+dY(t)
\end{aligned}\tag{11.2}
\end{equation}
旋转坐标。

支配物理过程的方程仍与式 (11.1) 相同，但写在新坐标系中则是
\begin{equation}
\begin{aligned}
X'(t) &= \alpha'X'(t-1)+\beta'Y'(t-1)+\varepsilon'(t)\\
Y'(t) &= \gamma'X'(t-1)+\delta'Y'(t-1)+\eta'(t).
\end{aligned}\tag{11.3}
\end{equation}
带撇的系数可由原来的（不带撇的）系数经矩阵乘法得到。同样有：
\[
\varepsilon'(t)=a\varepsilon(t)+b\eta(t),\qquad \eta'(t)=c\varepsilon(t)+d\eta(t).
\]
由于 $\varepsilon(t)$ 与 $\eta(t)$ 不相关，$\varepsilon'(t)$ 与 $\eta'(t)$ 将相关，于是我们不再具备“当前状态的分量 $X'(t)$ 与 $Y'(t)$ 在给定前一状态的分量 $X'(t-1)$ 与 $Y'(t-1)$ 时不相关”这一条件。因此，统计时间（如果存在的话）不再沿物理时间行进。

现在需要表明：我们可以这样选择参数 $a$、$b$、$c$、$d$，使统计时间与物理时间反向而行，即：使当前状态的分量 $X'(t)$ 与 $Y'(t)$ 在给定未来状态的分量 $X'(t+1)$ 与 $Y'(t+1)$ 时不相关。把式 (11.3) 反解，我们可以把 $X'(t-1)$ 与 $Y'(t-1)$ 表示为 $X'(t)$、$Y'(t)$、$\varepsilon'(t)$、$\eta'(t)$ 的线性组合。显然，由于 $\varepsilon(t)$ 与 $\eta(t)$ 不相关，我们可以这样选择 $a,b,c,d$，使出现在 $X'(t-1)$ 方程中的噪声项与出现在 $Y'(t-1)$ 方程中的噪声项不相关。（用矩阵演算可以更好地说明这一点。）

于是，选择替代坐标系的一般原则是：沿相反方向对角化噪声相关矩阵。

我希望读者接受这一挑战，去检验“时间偏倚猜想”（第 59 页）：

\begin{quote}
“在大多数自然现象中，物理时间与至少一个统计时间重合。”
\end{quote}

Alex Balke（私人通信）曾用经济时间序列试图检验它，但由于缺乏足够的数据，结果并不很有定论。我仍然相信该猜想为真，并希望读者运气更好。

\section{估计因果效应}

\subsection{后门准则背后的直觉（第三章，第 79 页）}

\qhead{问题（致作者）}

在后门条件的定义中（第 79 页，定义 3.3.1），对 $X$ 之后代的排除（条件 (i)）似乎是事后追加的，只是因为不这样做就会陷入麻烦。为什么不能从第一原理得到它：先按“消除偏倚”这一目标定义 $Z$ 的充分性，然后证明，为实现该目标，我们既不想要也不需要 $X$ 的后代进入 $Z$。

\qhead{作者答复}

把后代排除在后门准则之外，确实基于第一原理，其依据正是消除偏倚这一目标。原理如下：我们想测量某个量（因果效应），但实际测量到的却是由图中所有路径造成的相依 $P(y\mid x)$；其中有些是虚假的（后门路径），有些是真正因果的（从 $X$ 到 $Y$ 的有向路径）。因此，为消除偏倚，我们需要修改所测得的相依，使之等于所想要的量。为系统地做到这一点，我们对变量集 $Z$ 取条件，同时确保：

\begin{enumerate}[label=\arabic*.]
\item 阻断从 $X$ 到 $Y$ 的一切虚假路径；
\item 不扰动任何有向路径；
\item 不制造新的虚假路径。
\end{enumerate}

原理 1 与 2 通过阻断一切后门路径、且只阻断这些路径来实现，正如条件 (ii) 所表述。原理 3 要求我们不对 $X$ 的后代取条件，即便是那些不阻断有向路径的后代，因为这样的后代可能在 $X$ 与 $Y$ 之间制造新的虚假路径。要看清这一点，考虑图
\[
X\to S_1\to S_2\to S_3\to Y.
\]
中间变量 $S_1,S_2,\dots$（以及 $Y$）受噪声因子 $e_0,e_1,e_2,\dots$ 影响，这些因子未在图中显式画出。然而在放大之下，这条链展开为图 11.4 中的图。

现在设想我们对 $S_1$ 的一个后代 $Z$ 取条件，如图 11.5 所示。由于 $S_1$ 是碰撞点，这就在 $X$ 与 $e_1$ 之间造成了相依，其作用犹如一条后门路径
\[
X\leftrightarrow e_1\to S_1\to S_2\to S_3\to Y.
\]
按原理 3，这样的路径不应被制造出来，因为它们在 $X$ 与 $Y$ 之间引入了虚假相依。

注意，$X$ 的后代 $Z$ 若不同时是某个 $S_i$（或某个 $e_i$，$i>0$）的后代，则不受此排除：对它可以安全地取条件而不引入偏倚（尽管它可能降低 $X$ 对 $Y$ 之因果效应的相应估计量的效率）。\S 11.3.3 给出后门准则的另一种证明，其中排除 $X$ 之后代的必要性更加一目了然。

\begin{figure}[htbp]
\centering
\begin{tikzpicture}[causal,scale=0.95,every node/.style={transform shape}]
  \foreach \name/\x in {X/0,S1/1.5,S2/3.0,S3/4.5,Y/6.0}{%
    \node[circle,fill,inner sep=1.3pt] (\name) at (\x,0) {};
    \node[font=\small] at (\x,-0.34) {$\name$};
  }
  \node[font=\small] (e0) at (0,1.15) {$e_0$};
  \node[font=\small] (e1) at (1.5,1.15) {$e_1$};
  \node[font=\small] (e2) at (3.0,1.15) {$e_2$};
  \node[font=\small] (e3) at (4.5,1.15) {$e_3$};
  \node[font=\small] (e4) at (6.0,1.15) {$e_4$};
  \draw[ed] (X)--(S1); \draw[ed] (S1)--(S2); \draw[ed] (S2)--(S3); \draw[ed] (S3)--(Y);
  \draw[ed] (e0)--(X); \draw[ed] (e1)--(S1); \draw[ed] (e2)--(S2); \draw[ed] (e3)--(S3); \draw[ed] (e4)--(Y);
\end{tikzpicture}
\caption{显示从 $X$ 到 $Y$ 的路径上的噪声因子。}
\label{fig:11.4}
\end{figure}

\begin{figure}[htbp]
\centering
\begin{tikzpicture}[causal,scale=0.95,every node/.style={transform shape}]
  \foreach \name/\x in {X/0,S1/1.5,S2/3.0,S3/4.5,Y/6.0}{%
    \node[circle,fill,inner sep=1.3pt] (\name) at (\x,0) {};
    \node[font=\small] at (\x,-0.34) {$\name$};
  }
  \node[font=\small] (e0) at (0,1.15) {$e_0$};
  \node[font=\small] (e1) at (1.5,1.15) {$e_1$};
  \node[font=\small] (e2) at (3.0,1.15) {$e_2$};
  \node[font=\small] (e3) at (4.5,1.15) {$e_3$};
  \node[font=\small] (e4) at (6.0,1.15) {$e_4$};
  \coordinate (Zp) at (2.75,-1.55);
  \node[font=\small] at (2.88,-1.45) {$Z$};
  \draw[ed] (X)--(S1); \draw[ed] (S1)--(S2); \draw[ed] (S2)--(S3); \draw[ed] (S3)--(Y);
  \draw[ed] (e0)--(X); \draw[ed] (e1)--(S1); \draw[ed] (e2)--(S2); \draw[ed] (e3)--(S3); \draw[ed] (e4)--(Y);
  \draw[ed] (1.62,-0.28)--(Zp);
\end{tikzpicture}
\caption{对 $Z$ 取条件在 $X$ 与 $e_1$ 之间造成相依，从而使 $X$ 对 $Y$ 之效应的估计产生偏倚。}
\label{fig:11.5}
\end{figure}

同样重要的是要注意：为错误的变量作调整而制造新偏倚的危险，也会威胁随机化试验。在这类试验中，尽管随机化在渐近意义上同时中和了被测量与未被测量的混杂因子，研究者仍可能希望对协变量作调整。调整的动机，或者是提高精度（Cox 1958, pp.~48--55），或者是匹配不平衡的样本，或者是获得特定协变量下的因果效应。当调整仅限于处理前协变量时，随机化试验不会受调整所致偏倚的影响；但对处理后变量作调整，则可能通过图 11.5 所示的机制引入偏倚，或者更严重地，在调整变量 $Z$ 与影响结局的某个因子（如图 11.5 中的 $e_4$）之间存在相关时引入偏倚。

举例来说，假设处理对易患疾病 $Y$ 的病人有一种副作用（例如头痛）。如果我们想对“易感性”作调整，却转而对其代理变量“头痛”作调整，就会通过虚假路径——处理 $\to$ 头痛 $\leftrightarrow$ 易感性 $\to$ 疾病——产生偏倚。然而，只要我们谨慎地从不对处理之后果中那些会打开此类虚假路径的变量作调整（不只是位于通往疾病之因果通路上的那些），随机化试验中就不会出现偏倚。

\qhead{这位读者的进一步问题}

对排除 $X$ 之后代的这一解释是合理的，但有两个短处：

\begin{enumerate}[label=\arabic*.]
\item 它没有处理诸如
\[
X\leftarrow C\to Y\leftarrow F
\]
这样的情形；这类情形在流行病学中很常见，而传统允许对 $Z=\{C,F\}$ 作调整。
\item 这一解释似乎重新定义了混杂与充分性，使之表示的含义有别于过去几十年流行病学家所理解的含义。我们能否在图论中找到更接近其传统含义的东西？
\end{enumerate}

\qhead{作者答复}

\begin{enumerate}[label=\arabic*.]
\item 流行病学传统允许对 $Z=(C,F)$ 作调整，以完成“检验 $X$ 对 $Y$ 是否有因果效应”这一任务，但不允许用它估计效应的量值。就前一任务而言，虽然对 $F$ 取条件会在 $C$ 与影响 $Y$ 的噪声因子之间造出一条虚假路径，但该路径在对 $C$ 取条件时即被阻断。因此，对 $Z=\{C,F\}$ 取条件使 $X$ 与 $Y$ 保持独立。倘若我们在 $Z$ 的任何一层中恰好测到这种相依，那必定是模型有误——或者 $X$ 对 $Y$ 有直接因果效应，或者存在图中未画出的其他路径。

因此，如果我们想检验“$X$ 对 $Y$ 无因果效应”这一（零）假设，对 $Z=\{C,F\}$ 作调整是完全正当的，图也表明了这一点（即 $C$ 与 $F$ 都不是 $X$ 的后代）。然而，当怀疑 $X$ 对 $Y$ 有因果效应、因而要评估该效应时，对 $Z$ 作调整就不正当了，因为适用于这一任务的图是图 11.6；$F$ 成为 $X$ 的后代，被后门准则排除在外。

\item 如果对混杂与充分性的解释听起来与传统流行病学不一致，那只是因为传统流行病学家没有恰当的手段来表达“阻断”或“制造”相依这类操作。他们对相依可能已有健全的直觉，而图把这种直觉转化为一套关于路径开合的正式系统。

我们还应注意，1985 年以前，流行病学中的因果分析处于混乱状态，因为领军流行病学家的健全直觉不得不用关联的语言来表达——这是一项不可能完成的任务。甚至连“混杂代表偏倚”，即“两个相依之间的差别，一个是我们想测量的，一个是我们实际测量的”这一想法，也遭到许多人的抵制（见第六章），因为他们无法把前者用数学表达出来。\footnote{回想一下，Greenland and Robins (1986) 多年来是唯一一座真理的灯塔；而就连他们也只得借助“可交换性”这种“黑箱”语言来定义“偏倚”，这就妨碍了对混杂的直观解释（见 \S 6.5.3）。确实，流行病学家又花了六年（Weinberg 1993）才发现：对受暴露影响的因子（如图 11.5）作调整会引入偏倚。}

因此，与其在图的语言中寻找“更接近传统含义的东西”，我们不如做得更好：把那个“传统含义”本应是什么，明确地阐发出来。

换言之，传统含义是非正式的，而且偶尔会误入歧途；图形准则则是正式的、经过数学证明的。第六章（第 183、194 页）记录了一长串流行病学直觉的历史，其中一些出自著名流行病学家，却在混杂与调整问题上走入歧途（见 Greenland and Robins 1986; Wickramaratne and Holford 1987; Weinberg 1993）。尽管今天多数领军流行病学家对因果分析的现代发展极为敏感（例如 Glymour and Greenland 2008），流行病学的民俗中仍充斥着极为可疑的传统直觉。（见 \S 6.5.2。）

总之，图形准则以及上面的原理 1—3，为“流行病学概念的传统含义”提供了一种合理、友好而无歧义的解读。
\end{enumerate}

\begin{figure}[htbp]
\centering
\begin{tikzpicture}[causal,scale=0.95,every node/.style={transform shape}]
  \node[circle,fill,inner sep=1.3pt] (X) at (0,0) {};   \node[font=\small,below=0.5mm] at (0,0) {$X$};
  \node[circle,fill,inner sep=1.3pt] (Y) at (1.9,0) {}; \node[font=\small,below=0.5mm] at (1.9,0) {$Y$};
  \node[circle,fill,inner sep=1.3pt] (F) at (3.7,0) {}; \node[font=\small,below=0.5mm] at (3.7,0) {$F$};
  \node[circle,fill,inner sep=1.3pt] (C) at (1.0,1.15) {}; \node[font=\small,above=0.5mm] at (1.0,1.15) {$C$};
  \draw[ed] (X)--(Y); \draw[ed] (Y)--(F);
  \draw[ed] (C)--(X); \draw[ed] (C)--(Y);
\end{tikzpicture}
\caption{适用于评估 $X$ 对 $Y$ 之效应的图。}
\label{fig:11.6}
\end{figure}

\subsection{祛魅“强可忽略性”}

在潜在结局语言范围内工作的研究者，用名为“强可忽略性”的独立关系来表述“零偏倚”或“无混杂”条件（Rosenbaum and Rubin 1983）。形式上，若 $X$ 为二元处理（或行动），强可忽略性写作：
\begin{equation}
\{Y(0),Y(1)\}\indep X\mid Z, \tag{11.4}
\end{equation}
其中 $Y(0)$ 与 $Y(1)$ 分别是在行动 $\dox{X=0}$ 与 $\dox{X=1}$ 下的（不可观测的）潜在结局（定义见式 (3.51)），$Z$ 是一组被测量的协变量。当“强可忽略性”成立时，$Z$ 是可容许的，或曰去混杂的；也就是说，可以用调整估计量无偏地估计处理效应，如式 (3.54) 的推导所示。

正如该推导所示，强可忽略性是操纵反事实公式的一种便利的句法工具，也是形式上假定（$Z$ 的）可容许性而无须为之辩护的一种便利方式。然而，正如本书已多次指出的，几乎没有谁知道如何在实际中应用它，因为反事实变量 $Y(0)$ 与 $Y(1)$ 不可观测，而科学知识又不是以某种能对反事实的条件独立性作出可靠判断的形式储存的。因此毫不奇怪，“强可忽略性”几乎只被用作“$Z$ 可容许”这一假设的替身，即
\begin{equation}
P(y\mid\dox{x})=\sum_{z}P(y\mid z,x)P(z), \tag{11.5}
\end{equation}
而极少（如果有过的话）被用作保护我们免于糟糕 $Z$ 选择的准则。\footnote{事实上，在少数用“强可忽略性”指导协变量选择的情形中，所发布的准则都是错误或不准确的，从而使一些神话得以流传，例如：“没有理由避免对描述受试者处理前状况的变量作调整”；“混杂因子就是任何同时与处理和疾病相关的变量”；“强可忽略性要求测量一切与处理和结局都相关的协变量”（为免尴尬，此处不列出处）。}

受图形模型启迪的读者会立刻认识到，式 (11.4) 必定映射后门准则（第 79 页，定义 3.3.1），因为后者同样蕴含可容许性。这一认识使我们不仅能把式 (11.4) 当作一个主张或假设来提出，还能对使之成立的因果关系进行推理。

然而问题在于：能否把变量 $Y(0)$ 与 $Y(1)$ 以某种方式表示在因果图中，使我们能借助 d-分离、用图形手段检验式 (11.4)？换言之，我们要找一个节点集 $W$，使得 $Z$ d-分离 $X$ 与 $W$，当且仅当 $Z$ 满足式 (11.4)。

答案直接来自图与潜在结局之间的翻译规则（\S 3.6.3）。按此翻译，$\{Y(0),Y(1)\}$ 表示一切外生变量——无论潜变量还是可观测变量——的总和，这些变量能经由避开 $X$ 的路径影响 $Y$。理由如下：按 $\{Y(0),Y(1)\}$ 的结构定义（式 (3.51)），$Y(0)$（$Y(1)$ 同理）表示在“一切进入 $X$ 的箭头都被切断、且 $X$ 被保持在 $X=0$”这一条件下 $Y$ 的取值。因此，$Y(0)$ 的统计变异，受控于切除图中（切去进入 $X$ 的箭头后）$Y$ 的一切外生祖先。

例如在图 11.4 中，$\{Y(0),Y(1)\}$ 将由外生变量 $\{e_1,e_2,e_3,e_4\}$ 表示。再如，在图 3.4 中，$\{Y(0),Y(1)\}$ 将由影响变量 $X_4$、$X_1$、$X_2$、$X_5$、$X_6$ 的噪声因子（图中未画出）表示。然而，由于变量 $X_4$ 与 $X_5$ 对（$Y$ 而言）概括了其祖先的变异，表示 $\{Y(0),Y(1)\}$ 的一个充分集合将是 $X_4$、$X_1$ 以及影响 $Y$ 与 $X_6$ 的噪声因子。

总之，潜在结局 $\{Y(0),Y(1)\}$ 由从 $X$ 到 $Y$ 的路径上所有节点的、被观测与未被观测的父节点\footnote{把潜父节点显式包含进来的理由，见 \S 11.3.1。}表示。示意地，我们可以把这些父节点表示为图 11.7(a)。容易看出，在对 $\{Y(0),Y(1)\}$ 作此解释之下，协变量集 $Z$ d-分离 $W$ 与 $X$，当且仅当 $Z$ 满足后门准则。

应当注意，图 11.7(a) 中被标为 $W$ 的那组可观测变量，只是为了在因果图中（经由 d-分离）确认条件独立（例如式 (11.4)），才充当不可观测反事实 $\{Y(0),Y(1)\}$ 的代理。$\{Y(0),Y(1)\}$ 更精确的配置见图 11.7(b)：那里它们被画成 $Y$ 的（哑）父节点，是 $Y$ 与 $S$ 的实际（可观测）父节点的函数，但并不等同于后者。

\begin{figure}[htbp]
\centering
\begin{tikzpicture}[causal,scale=0.9,every node/.style={transform shape}]
  % (a)
  \begin{scope}
    \node[circle,fill,inner sep=1.3pt] (aX) at (0,0) {};   \node[font=\small,below=0.7mm] at (0,0) {$X$};
    \node[circle,fill,inner sep=1.3pt] (aS) at (2.1,0) {}; \node[font=\small,below=0.7mm] at (2.1,0) {$S$};
    \node[circle,fill,inner sep=1.3pt] (aY) at (4.3,0) {}; \node[font=\small,below=0.7mm] at (4.3,0) {$Y$};
    \draw[ed] (aX)--(aS); \draw[ed] (aS)--(aY);
    \draw[dashed] (2.85,1.0) ellipse (1.75cm and 0.62cm);
    \node[circle,fill,inner sep=1.2pt] (ap0) at (0.35,1.4) {};
    \node[circle,fill,inner sep=1.2pt] (ap1) at (1.3,0.7) {};
    \node[circle,fill,inner sep=1.2pt] (ap2) at (3.05,1.4) {};
    \node[circle,fill,inner sep=1.2pt] (ap3) at (3.9,0.7) {};
    \draw[ed] (ap0)--(aX); \draw[ed] (ap0)--(ap1); \draw[ed] (ap1)--(aS);
    \draw[ed] (ap2)--(aS); \draw[ed] (ap2)--(ap3); \draw[ed] (ap3)--(aY);
    \node[font=\small] (aW) at (5.25,1.75) {$W$};
    \draw[->,dashed] (aW) -- (4.45,1.25);
    \node[font=\small] at (2.15,-0.9) {(a)};
  \end{scope}
  % (b)
  \begin{scope}[xshift=7.5cm]
    \node[circle,fill,inner sep=1.3pt] (bX) at (0,0) {};   \node[font=\small,below=0.7mm] at (0,0) {$X$};
    \node[circle,fill,inner sep=1.3pt] (bS) at (2.1,0) {}; \node[font=\small,below=0.7mm] at (2.1,0) {$S$};
    \node[circle,fill,inner sep=1.3pt] (bY) at (4.3,0) {}; \node[font=\small,below=0.7mm] at (4.3,0) {$Y$};
    \draw[ed] (bX)--(bS); \draw[ed] (bS)--(bY);
    \node[circle,fill,inner sep=1.2pt] (bp0) at (0.35,1.4) {};
    \node[circle,fill,inner sep=1.2pt] (bp1) at (1.3,0.7) {};
    \node[circle,fill,inner sep=1.2pt] (bp2) at (3.05,1.4) {};
    \node[circle,fill,inner sep=1.2pt] (bp3) at (3.9,0.7) {};
    \draw[ed] (bp0)--(bX); \draw[ed] (bp0)--(bp1); \draw[ed] (bp1)--(bS);
    \draw[ed] (bp2)--(bS); \draw[ed] (bp2)--(bp3); \draw[ed] (bp3)--(bY);
    \draw[thick] (5.35,1.45) ellipse (0.62cm and 0.30cm);
    \node[font=\small,right=0.8mm] at (5.98,1.45) {$\{Y(0),Y(1)\}$};
    \draw[edd] (bp2)--(4.9,1.62);
    \draw[edd] (bp3)--(5.0,1.3);
    \draw[ed] (5.35,1.15)--(bY);
    \node[font=\small] at (2.15,-0.9) {(b)};
  \end{scope}
\end{tikzpicture}
\caption{“强可忽略性”条件下反事实 $\{Y(0),Y(1)\}$ 的图形解释：(a) 用可观测代理 $W$ 表示这些反事实所涵盖的父节点集合；(b) 把 $\{Y(0),Y(1)\}$ 本身画成 $Y$ 的哑父节点。（示意图。）}
\label{fig:11.7}
\end{figure}

熟悉结构方程建模的读者会认出，$\{Y(0),Y(1)\}$ 的图形表示是对计量经济学中“扰动”或（$Y$ 方程中的）“误差项”这一经典概念的精细化；而“强可忽略性”则是这样的要求：$X$ 若要“外生”，就必须与这一“扰动”独立（见 \S 5.4.3）。这一概念在 1970 年代与计量经济方程的因果解释一起声名扫地（Richard 1980）；鉴于图形模型为结构方程形式体系带来的清明，我曾预言它会被重新接受（第 170 页）。图 11.7 应会进一步助成这种接受。

把“强可忽略性”翻译成一个编码了实质过程知识的模型中的简单分离条件，应当能祛除“强可忽略性”这一含混概念的神秘性，并吸引讲“可忽略性”语言的研究者从它的图形解释中获益。

这一解释使研究者能够理解：协变量在消除偏倚之前必须满足什么条件；在选择协变量时应注意什么、考虑什么；以及我们可以做哪些实验，来至少部分地检验我们是否具备协变量选择所需的知识。\S 11.3.4 举例说明了这类考量。

图形框架与反事实框架之间的共生关系曾在一个应用中大显身手，即估计处理对受处理者的效应：$\mathrm{ETT}=P(Y_{x'}=y\mid x)$（见 \S 8.2.5 与 \S 11.9.1）。这一反事实量（例如：受处理者若不受处理将会康复的概率，或者暴露者在避免暴露时的疾病率）在 do 演算记号中不易分析。然而反事实记号使我们能推出一个有用的结论：只要存在满足后门准则的协变量集 $Z$，ETT 就可以从观测研究中估计。这直接来自
\[
(Y\indep X\mid Z)_{G_{\underline{X}}}\ \Longrightarrow\ Y_{x'}\indep X\mid Z,
\]
由此可以写出
\[
\begin{aligned}
\mathrm{ETT} &= P(Y_{x'}=y\mid x)\\
&= \sum_{z}P(Y_{x'}=y\mid x,z)P(z\mid x)\\
&= \sum_{z}P(Y_{x'}=y\mid x',z)P(z\mid x)\\
&= \sum_{z}P(y\mid x',z)P(z\mid x).
\end{aligned}
\]
“强可忽略性”的图形祛魅还有助于解释：为什么因果概率 $P(Y_{x'}=y'\mid x,y)$，事实上还有任何以 $y$ 为条件的反事实表达式，都不允许这样的推导，并且一般是不可识别的（见第九章与 Shpitser and Pearl 2007）。

\subsection{后门准则的另一证明}

后门准则（定理 3.3.2）的原始证明借助了一个辅助干预节点 $F$（图 3.2），相当间接。下面给出另一种证明，其中把 $Z$ 限制为 $X$ 之非后代的必要性一目了然。

\subsubsection*{后门准则的证明}

考虑一个马尔可夫模型 $G$，其中 $T$ 表示 $X$ 的父节点集。由式 (3.13) 可知，$X$ 对 $Y$ 的因果效应为
\begin{equation}
P(y\mid\hat{x})=\sum_{t\in T}P(y\mid x,t)P(t). \tag{11.6}
\end{equation}
现在假定 $T$ 中某些成员未被观测。我们要另找一组可观测变量 $Z$ 来替代 $T$，使得
\begin{equation}
P(y\mid\hat{x})=\sum_{z\in Z}P(y\mid x,z)P(z). \tag{11.7}
\end{equation}
容易验证：若 $Z$ 满足
\begin{enumerate}[label=(\roman*)]
\item $(Y\indep T\mid X,Z)$
\item $(X\indep Z\mid T)$，
\end{enumerate}
则 (11.7) 由 (11.6) 推出。确实，对 $Z$ 取条件，(i) 使我们能把 (11.6) 改写为
\[
P(y\mid\hat{x})=\sum_{t}P(t)\sum_{z}P(y\mid z,x)P(z\mid t,x),
\]
而 (ii) 进一步给出 $P(z\mid t,x)=P(z\mid t)$，于是 (11.7) 得证。

现在证明“后门准则蕴含 (i) 与 (ii)”就纯属图形练习了。确实，(ii) 直接来自 $Z$ 由 $X$ 的非后代构成这一事实；而 $Z$ 阻断一切后门路径则蕴含 $(Y\indep T\mid X,Z)_G$，从而 (i) 成立。后者来自如下观察：$G$ 中任何从 $Y$ 到 $T$、未被 $\{X,Z\}$ 阻断的路径，都可延拓为一条从 $Y$ 到 $X$、未被 $Z$ 阻断的后门路径。

\subsubsection*{关于可容许去混杂集的识别}

注意，条件 (i) 与 (ii) 使我们能从任何其他可容许的 $T$ 出发，识别出一个集合 $Z$ 是否可容许（即是否满足式 (11.7)），而 $T$ 不必是 $X$ 的父节点。$T$ 的父节点身份只用于建立 (11.6)，在“以 $Z$ 替换 $T$ 以建立 (11.7)”中不起作用。尽管如此，从父节点集 $T$ 出发有一个独特的优势：它使我们能经由 (i) 与 (ii) 识别出其他每一个可容许集 $Z$。对任何别的起点 $T$，都存在一个不满足 (i) 与 (ii) 的可容许 $Z$。一个明显的例子是：取 $X$ 的父节点为 $Z$ 会违反 (i) 与 (ii)，因为如 (i) 所要求，没有任何集合能把 $X$ 与其父节点 d-分离。

还要注意，(i) 与 (ii) 是纯统计的，不援引关于图的知识或任何其他因果假设。因此有趣的是：是否存在与 (i)、(ii) 类似的一般独立条件，把任意两个可容许集 $S_1$ 与 $S_2$ 联系起来？部分答案由 Stone–Robins 准则（第 187 页）给出，它处理 $S_1$ 是 $S_2$ 之子集的情形；另一部分答案由下述观察给出。

若下述等式成立，则称两个子集 $S_1$ 与 $S_2$ 是 c-等价的（“c”表示混杂）：
\begin{equation}
\sum_{s_1}P(y\mid x,s_1)P(s_1)=\sum_{s_2}P(y\mid x,s_2)P(s_2). \tag{11.8}
\end{equation}
这一等式保证：若对之调整，$S_1$ 与 $S_2$ 在估计 $X$ 对 $Y$ 的因果效应方面会产生相同的偏倚。

\said{断言：}$S_1$ 与 $S_2$ c-等价的一个充分条件，是下述两个条件之一成立：
\[
\begin{aligned}
C1:\quad & X\indep S_2\mid S_1 &&\text{且}\quad Y\indep S_1\mid S_2,X\\
C2:\quad & X\indep S_1\mid S_2 &&\text{且}\quad Y\indep S_2\mid S_1,X
\end{aligned}
\]
C1 使我们能从式 (11.8) 的左边推出右边，C2 则使我们能从右边推回左边。因此，若已知 $S_1$ 可容许，则 $S_2$ 的可容许性可由 C1 或 C2 确认。这一更宽的条件使我们能够，例如，从任何其他可容许集 $S_1$ 出发证明 $S_2=\PA_X$ 可容许，因为任何这样的选择都满足条件 C2。

这一更宽的条件仍未能刻画所有 c-等价对 $S_1$、$S_2$。例如考虑图 11.8(a)，其中 $S_1=\{Z_1,W_2\}$ 与 $S_2=\{Z_2,W_1\}$ 各自可容许（因为它们满足后门准则），故 $S_1$ 与 $S_2$ 是 c-等价的。然而在此情形下，C1 与 C2 都不成立。

\begin{figure}[htbp]
\centering
\begin{tikzpicture}[causal,scale=0.9,every node/.style={transform shape}]
  % (a)
  \begin{scope}
    \node[circle,fill,inner sep=1.3pt] (aW1) at (0,2.3) {};   \node[font=\small,above=0.6mm] at (0,2.3) {$W_1$};
    \node[circle,fill,inner sep=1.3pt] (aW2) at (2.7,2.3) {}; \node[font=\small,above=0.6mm] at (2.7,2.3) {$W_2$};
    \node[circle,fill,inner sep=1.3pt] (aZ1) at (0.55,1.2) {}; \node[font=\small,below right=0.1mm and 0.6mm] at (0.55,1.2) {$Z_1$};
    \node[circle,fill,inner sep=1.3pt] (aZ2) at (2.15,1.2) {}; \node[font=\small,below left=0.1mm and 0.6mm] at (2.15,1.2) {$Z_2$};
    \node[circle,fill,inner sep=1.3pt] (aX) at (0,0) {};       \node[font=\small,below=0.7mm] at (0,0) {$X$};
    \node[circle,fill,inner sep=1.3pt] (aY) at (2.7,0) {};     \node[font=\small,below=0.7mm] at (2.7,0) {$Y$};
    \draw[ed] (aW1)--(aW2); \draw[ed] (aZ1)--(aZ2);
    \draw[ed] (aW1)--(aX); \draw[ed] (aZ1)--(aX);
    \draw[ed] (aW2)--(aY); \draw[ed] (aZ2)--(aY);
    \draw[ed] (aX)--(aY);
    \node[font=\small] at (1.35,-0.95) {(a)};
  \end{scope}
  % (b)
  \begin{scope}[xshift=6.9cm]
    \node[circle,fill,inner sep=1.3pt] (bW1) at (0,2.3) {};   \node[font=\small,above=0.6mm] at (0,2.3) {$W_1$};
    \node[circle,fill,inner sep=1.3pt] (bW2) at (2.7,2.3) {}; \node[font=\small,above=0.6mm] at (2.7,2.3) {$W_2$};
    \node[circle,fill,inner sep=1.3pt] (bZ1) at (0.55,1.2) {}; \node[font=\small,below right=0.1mm and 0.6mm] at (0.55,1.2) {$Z_1$};
    \node[circle,fill,inner sep=1.3pt] (bZ2) at (2.15,1.2) {}; \node[font=\small,below left=0.1mm and 0.6mm] at (2.15,1.2) {$Z_2$};
    \node[circle,fill,inner sep=1.3pt] (bV) at (1.35,1.75) {};  \node[font=\small,above=0.6mm] at (1.35,1.75) {$V$};
    \node[circle,fill,inner sep=1.3pt] (bX) at (0,0) {};       \node[font=\small,below=0.7mm] at (0,0) {$X$};
    \node[circle,fill,inner sep=1.3pt] (bY) at (2.7,0) {};     \node[font=\small,below=0.7mm] at (2.7,0) {$Y$};
    \draw[dashed] (1.35,1.5) ellipse (1.9cm and 1.28cm);
    \draw[ed] (bW1)--(bZ1); \draw[ed] (bW2)--(bZ2);
    \draw[ed] (bV)--(bZ1); \draw[ed] (bV)--(bZ2);
    \draw[ed] (bZ1)--(bX); \draw[ed] (bZ2)--(bY);
    \draw[ed] (bX)--(bY);
    \node[font=\small] at (1.35,-0.95) {(b)};
  \end{scope}
\end{tikzpicture}
\caption{(a) $S_1=\{Z_1,W_2\}$ 与 $S_2=\{Z_2,W_1\}$ 各自可容许，却不满足 $C_1$ 或 $C_2$。(b) $C=\{Z_1,Z_2,W_1,W_2,V\}$ 的任何子集都不可容许（虚线圆表示圈内变量共受一个未观测的共同原因影响）。}
\label{fig:11.8}
\end{figure}

一个自然的尝试是加上“$S_1$ 与 $S_2$ 各自与 $S_1\cup S_2$ c-等价”这一条件，并援引 Stone (1993) 与 Robins (1997) 关于所需集合—子集等价的准则。所得准则虽然有效，仍不完备：有些情形下 $S_1$ 与 $S_2$ c-等价，却都不与其并集 c-等价。Pearl and Paz (2008) 的一个定理用不可约集拓宽了这一条件。

当然，给出 c-等价的条件独立刻画，并没有解决识别可容许集的问题；后者是因果概念，无法给出统计刻画。

图 11.8(b) 所示的图，展示了社会科学家与健康科学家常面临的困难。假设我们的目标是在测得 $\{X,Y,Z_1,Z_2,W_1,W_2,V\}$ 的条件下估计 $P(y\mid\dox{x})$，却对底层的图结构一无所知。常规做法是从一切可得的协变量 $C=\{Z_1,Z_2,W_1,W_2,V\}$ 出发，检验 $C$ 的某个真子集在调整后是否会给出等价的估计量。这类约简的统计方法见 Greenland et al. (1999b)、Geng et al. (2002) 与 Wang et al. (2008)。例如，相继施用条件 C1 与 C2，可以把 $V$ 与 $Z_2$ 从 $C$ 中删去，从而得到一个与原协变量集 $C$ c-等价的不可约子集 $\{Z_1,W_1,W_2\}$。然而这个子集不可用于调整，因为它与 $C$ 一样不满足后门准则。Tian et al. (1998) 的一个定理向我们保证：$C$ 的任何 c-等价子集都可由 $C$ 经“一次删一个”的方法、经过一串 c-等价子集到达；但协变量选择的问题在于：由于缺乏图结构，我们不知道 $C$ 的众多子集中哪一个（如果有的话）可容许。下一小节讨论如何借助外部知识以及对现有数据的更精细分析来处理这个问题。

\subsection{协变量选择中的数据与知识}

那么，在没有因果图的情况下能做什么？一种办法是：基于对领域的理解，假定一个看似合理的图，然后检验数据是否否证该图所蕴含的任何统计主张。在我们的例子里，图 11.8(b) 的图宣称了若干此类主张，它们以条件独立约束的形式出现，每一条都对应于图中一条缺失的箭头：
\[
\begin{aligned}
V &\indep \{W_1,W_2\} && X\indep\{V,Z_2\}\mid\{Z_1,W_2,W_1\}\\
Z_1 &\indep \{W_2,Z_2\}\mid\{V,W_1\} && V\indep Y\mid\{X,Z_2,W_2,Z_1,W_1\}\\
Z_2 &\indep \{W_1,Z_1,X\}\mid\{V,W_2\} && V\indep Y\mid\{Z_2,W_2,Z_1,W_1\}
\end{aligned}
\]

满足这些约束当然不能确立所设定因果模型的有效性，因为如第二章所示，可能存在满足同样独立约束、却体现显著不同因果结构的替代模型，因而具有显著不同的可容许集与效应估计量。一个平凡的例子是完全图：其箭头方向可以任意，只要巧妙选择参数，它就能模拟任何其他图。一个不那么平凡、且对参数选择不敏感的例子，是等价结构类：其中一切条件独立都来自图的分离。\footnote{半马尔可夫模型还可以由无法表述为条件独立的函数关系来区分（Verma and Pearl 1990; Tian and Pearl 2002b; Shpitser and Pearl 2008）。本例不把这些有用的约束考虑在内。}第二章所发展的搜索技术，提供了系统地表示与给定条件独立关系集相容的所有等价模型的方法。

图 11.9 所示的模型是图 11.8(b) 的强劲对手：它满足后者所蕴含的一切条件独立，还多一条
\[
Z_1\indep Y\mid X,W_1,W_2,Z_2,
\]
而这一条不易察觉或检验。然而与图 11.8(b) 相反，它认为三个集合 $\{Z_1,W_1,W_2\}$、$\{V,W_1,W_2\}$ 与 $\{Z_2,W_1,W_2\}$ 可容许，因而 c-等价；检验这三个集合的 c-等价性，应当能在两个竞争的模型之间作出裁决。

\begin{figure}[htbp]
\centering
\begin{tikzpicture}[causal,scale=0.85,every node/.style={transform shape}]
  \node[circle,fill,inner sep=1.4pt] (V) at (3.0,4.2) {};   \node[font=\small,above=0.6mm] at (3.0,4.2) {$V$};
  \node[circle,fill,inner sep=1.4pt] (W1) at (1.4,1.9) {};  \node[font=\small,above=0.7mm] at (1.4,1.9) {$W_1$};
  \node[circle,fill,inner sep=1.4pt] (W2) at (4.6,1.9) {};  \node[font=\small,above=0.7mm] at (4.6,1.9) {$W_2$};
  \node[circle,fill,inner sep=1.4pt] (Z1) at (-0.4,1.0) {}; \node[font=\small,left=0.7mm] at (-0.4,1.0) {$Z_1$};
  \node[circle,fill,inner sep=1.4pt] (Z2) at (6.4,1.0) {};  \node[font=\small,right=0.7mm] at (6.4,1.0) {$Z_2$};
  \node[circle,fill,inner sep=1.4pt] (X) at (1.0,0) {};     \node[font=\small,below=0.7mm] at (1.0,0) {$X$};
  \node[circle,fill,inner sep=1.4pt] (Y) at (5.0,0) {};     \node[font=\small,below=0.7mm] at (5.0,0) {$Y$};
  \draw[ed] (V)--(Z1); \draw[ed] (V)--(Z2);
  \draw[ed] (W1)--(Z1); \draw[ed] (W2)--(Z2);
  \draw[ed] (W1)--(X);  \draw[ed] (W1)--(Y);
  \draw[ed] (W2)--(X);  \draw[ed] (W2)--(Y);
  \draw[ed] (Z1)--(X);  \draw[ed] (Z2)--(Y);
  \draw[ed] (X)--(Y);
  \draw[bid,dashed] (W1) to[bend left=32] (W2);
\end{tikzpicture}
\caption{一个与图 11.8(b) 几乎无法区分的模型，只是多宣称了一条独立关系 $Z_1\indep Y\mid X,W_1,W_2,Z_2$。它认为三个集合 $\{Z_1,W_1,W_2\}$、$\{V,W_1,W_2\}$ 与 $\{Z_2,W_1,W_2\}$ 可容许（因而 c-等价）；因此，只要它们之中有任何一对未通过 c-等价检验，该模型就会被拒斥。}
\label{fig:11.9}
\end{figure}

实质性因果知识可以为这类裁决提供有价值的信息。例如，如果我们有充分理由相信变量 $W_2$ 不可能对 $X$ 有任何影响（例如它可能发生在 $X$ 之后），或者 $W_1$ 不可能对 $Y$ 有直接效应，那么图 11.9 的模型就可以被排除。

图的力量在于为研究者提供一种透明的语言，用来推理、讨论此类假设的合理性，并在不能达成共识时，隔离意见分歧、指明需要哪些额外观测才能解决分歧。潜在结局进路缺乏这种便利：对多数研究者而言，“强可忽略性”始终是一个神秘的黑箱。

除了充当实质判断的载体，图还使人能够在不逐一检验类中成员的情况下拒斥整类模型。例如，所有“以 $V$ 与 $W_1$ 为 $X$ 的唯一父节点、从而使 $\{V,W_1\}$（以及 $C$）可容许”的模型，只要数据中条件
\[
X\indep Z_1\mid V,W_1
\]
不成立，就可以一下子全部拒斥。

例如在第三章我们演示过：额外测量一个在 $X$ 与 $Y$ 之间起中介作用的变量，足以识别 $X$ 对 $Y$ 的因果效应。这一便利也可在图 11.8(b) 中演示：测量一个被判断位于 $X$ 与 $Y$ 之间通路上的变量 $Z$，将使 $P(y\mid\dox{x})$ 可识别，并可通过式 (3.29) 估计。当然，这一切以图 11.8(b) 为正确的数据生成模型为前提。反之，如果代表正确模型的是图 11.9，则因果效应将由
\[
P(y\mid\dox{x})=\sum_{\pa_X}P(y\mid\pa_X,x)P(\pa_X)
=\sum_{z_1,w_1,w_2}P(y\mid x,z_1,w_1,w_2)P(z_1,w_1,w_2)
\]
给出，这可能与式 (3.29) 一致，也可能不一致。在后一情形，我们就有充分理由认为图 11.9 不一致而拒斥它，或许还会寻求额外测量以确认或否证图 11.8(b)。

辅助实验可能提供比辅助观测更有力的判别工具。考虑图 11.8(b) 中的变量 $W_1$。如果我们能做一个以 $W_1$（而非 $X$）为随机化对象的受控实验，所得数据将使我们能无偏地估计 $X$ 对 $Y$ 的因果效应（见 \S 3.4.4）。至少，我们将能辨明 $W_1$ 究竟是图 11.9 那样为 $X$ 的父节点，还是图 11.8(b) 那样为 $X$ 的间接祖先。

为力图恪守传统统计方法，一些因果分析者采用了一种称为“敏感性分析”的方法（例如 Rosenbaum 2002, pp.~105--170），它给人的印象是分析中并未援引因果假设。这当然是一种错觉。分析者不是通过假定模型中不存在某些因果关系来作推断，而是尝试各种这样的假定，并评估替代的因果关系必须多强才能解释所观测到的数据。结果随后交由合理性判断，而这种判断的性质，与设定图 11.9 那样的模型时所援引的判断并无不同。在更丰富的设定中，敏感性分析相当于给一个图赋予若干因果关系、其强度受合理性判断限制，然后在给定数据下，试图在不违反这些合理性约束的前提下得出结论。这是一项高尚的努力，但迄今限于变量数很少的问题。图的出现有望把这一方法的适用范围扩展到更现实的问题。

\subsection{理解倾向评分}

倾向评分法（Rosenbaum and Rubin 1983），或称倾向评分匹配（PSM），是观测研究中发展最充分、最流行的因果分析策略。本书不强调它，因为它是估计方法，旨在处理有限样本的变异性，对本书主要关注的渐近、大样本极限的理解帮助不大。然而，由于倾向评分法在因果分析中的显赫地位以及近来围绕其使用的争议，我们专辟一节，说明它在图形模型、可容许性、可识别性、偏倚削减以及统计与因果二分这一大格局中的位置。

倾向评分法建立在一个简单却巧妙的、纯统计性质的想法上。假定行动（或处理）$X$ 为二元，$S$ 为任意的被测量协变量集，则倾向评分 $L(s)$ 是特征为 $S=s$ 的参与者选择行动 $X=1$ 的概率，即
\begin{equation}
L(s)=P(X=1\mid S=s). \tag{11.9}
\end{equation}
Rosenbaum 与 Rubin 所证明的是：把 $L(s)$ 视为 $S$ 的函数，因而视为随机变量，则给定 $L(s)$ 时 $X$ 与 $S$ 独立，即
\[
X\indep S\mid L(s).
\]
用文字说：所有映射到同一 $L(s)$ 值的单元是可比的，或曰“平衡的”，意思是：在 $L$ 的每一层内，受处理与未受处理单元具有相同的特征分布 $S$。\footnote{这一独立性源自函数 $L(s)$ 的特殊性质，并不在图中得到表示；也就是说，若把 $L$ 画成 $S$ 的子节点，$L$ 一般并不能把 $S$ 与 $X$ d-分离。}

为看清这一结果的意义，让我们为简单起见假定 $L(s)$ 可以单独估计，并可用离散层 $L=\{l_1,l_2,\dots,l_k\}$ 近似。条件独立 $X\indep S\mid L(s)$ 连同函数映射 $S\mapsto L$，使 $S$ 与 $L$ 在 \S 11.3.3 式 (11.8) 所定义的意义上 c-等价，即对任意 $Y$，
\begin{equation}
\sum_{s}P(y\mid s,x)P(s)=\sum_{l}P(y\mid l,x)P(l). \tag{11.10}
\end{equation}
这可由下式立即得出：\footnote{这也可由如下事实得出：把 $S_1=S$ 与 $S_2=L(s)$ 代入，即满足条件 C2。}
\[
\begin{aligned}
\sum_{l}P(y\mid l,x)P(l) &= \sum_{s}\sum_{l}P(y\mid l,s,x)P(l)P(s\mid l,x)\\
&= \sum_{s}\sum_{l}P(y\mid s,x)P(l)P(s\mid l)\\
&= \sum_{s}P(y\mid s,x)P(s).
\end{aligned}
\]

至此我们还没有提到任何因果关系，也没有提到 $Y$ 是结局变量、我们的任务终将是估计 $X$ 对 $Y$ 的因果效应。$S$ 与 $L$ 的 c-等价仅仅意味着：如果出于某种原因，人们想估计“调整估计量” $\sum_s P(y\mid s,x)P(s)$（其中 $S$ 与 $Y$ 是两个任意变量集），那么与其对高维集合 $S$ 求和，不如对一维向量 $L(s)$ 求和。在非常大样本的极限下，两种方法的渐近估计相同。

这一 c-等价进一步意味着——因果推断正是在这里首次登场——如果人们选择用调整估计量 $E_sP(y\mid s,x)$ 来近似因果效应 $P(y\mid\dox{x})$，那么在渐近意义上，用估计量 $E_lP(y\mid l,x)$ 也能达到同样的近似，其中调整是在 $L$ 的各层上进行的。后者有一个好处：对有限样本，每一层不太可能是空的，而且每一层都可能同时包含受处理与未受处理单元。

于是，倾向评分法可以看作调整估计量的一个有效估计量，而该调整估计量由任意的协变量集 $S$ 构成；它既不对 $S$ 是否恰当作任何陈述，也不承诺校正任何混杂偏倚，也不承诺在不存在偏倚的地方避免制造新偏倚。

在 $S$ 可容许的特殊情形，即
\begin{equation}
P(y\mid\dox{x})=E_sP(y\mid s,x), \tag{11.11}
\end{equation}
$L$ 也将可容许，于是我们就有了因果效应的一个无偏估计量\footnote{Rosenbaum 与 Rubin (1983) 只对可容许的 $S$ 证明了 $S$ 与 $L$ 的 c-等价，这很不幸：它使读者以为倾向评分匹配以某种方式有助于削减偏倚。}
\[
P(y\mid\dox{x})=E_lP(y\mid l,x),
\]
并附有一个估计右端的有效方法。反之，若 $S$ 不可容许，$L$ 也将不可容许，我们所能保证的只是：前者产生的偏倚会被后者忠实而有效地复制出来。

\subsubsection*{围绕倾向评分的争议}

至此，我们对倾向评分的陈述没有留下误解的余地，本书读者很难理解：一个无邪的估计方法，只不过提供了一个有效途径去估计某个统计量（该统计量有时与所关心的因果量重合、有时不重合，取决于 $S$ 的选择），怎么会引出争议呢？

但争议近来确实出现了，最可能的原因是这一方法日益流行，以及它得到了著名统计学家（Rubin 2007）、社会科学家（Morgan and Winship 2007; Berk and de Leeuw 1999）、健康科学家（Austin 2007）与经济学家（Heckman 1992）的强力背书。这一方法的流行程度甚至到了这样的地步：一些联邦机构如今期望项目评估者用这一方法替代实验设计（Peikes et al. 2008）。这一动向反映了一种普遍倾向：研究者淡化关于 $S$ 必须可容许的告诫，并把 Rosenbaum 与 Rubin 的数学证明解读为一项保证——在 $L$ 的每一层内，把受处理与未受处理受试者匹配起来，就能以某种方式消除数据中的混杂，从而对总体偏倚的削减有所贡献。这一倾向又被一些经验研究进一步强化（Heckman et al. 1998; Dehejia and Wahba 1999）：这些研究发现倾向评分分析与随机化试验结果一致，并把这种一致归因于前者能在重要特征上“平衡”处理组与对照组。Rubin 曾以如下陈述鼓励这类解读：“这一应用使用倾向评分方法创建受处理单元与对照单元的子组……就好像它们是被随机化的一样。这些子组的集合于是相对于被观测协变量‘近似’一个随机化区组实验”（Rubin 2007）。

然而，随后的经验研究对倾向评分采取了更具批判性的看法，失望地指出：当与临床研究的结果作仔细比较时，有时会测到可观的偏倚（Smith and Todd 2005; Luellen et al. 2005; Peikes et al. 2008）。

但是，为什么会有人淡化 Rosenbaum 与 Rubin 的告诫呢？这样做会违反因果分析的黄金法则：没有任何因果主张能由纯统计方法确立，无论是倾向评分、回归、分层，还是任何其他基于分布的设计。我认为，答案在于 Rosenbaum 与 Rubin 表述可容许条件（即式 (11.11)）时所使用的语言。该条件是用潜在结局的隐晦语言表述的：集合 $S$ 必须使 $X$ “强可忽略”，即
\[
\{Y_1,Y_0\}\indep X\mid S.
\]
正如本书多次指出的，“可忽略性”的晦涩是潜在结局进路的阿喀琉斯之踵——没有凡人能应用这一条件来判断它是否成立，即使在一切因果关系都被正确设定的简单问题中也是如此，更不用说涉及几十个变量的部分设定问题了。\footnote{请潜在结局传统的拥护者检视图 11.8(b)（或他们自选的任何模型、故事、玩具例子），判断 $C$ 的某个子集是否使 $X$“强可忽略”。当然，这用后门准则很容易判定；但不幸的是，图仍然被潜在结局阵营的一些主要拥护者所畏惧和误解（例如 Rubin 2004, 2008b, 2009）。}

多数研究者在理解“可忽略性”是什么意思、以及它召唤他们作何种判断时所体验到的困难，诱使他们假定：只要把尽可能多的协变量纳入分析，它就会自动得到满足，或者至少很可能得到满足。流行的态度是：加入更多协变量不会有害（Rosenbaum 2002, p.~76），并且可以免去人们去思考这些协变量、处理、结局之间——最重要的是未被测量的混杂因子之间——的因果关系（Rubin 2009）。

这一态度与图形模型的学生所学到的东西、以及本书试图传授的东西正相反。$S$ 的可容许性只能靠援引研究者可得的因果知识来确立；而正如我们从图论与后门准则所知，这种知识使偏倚削减成为非单调的操作：消除由某个混杂因子造成的偏倚（或不平衡），可能唤醒并释放由休眠的、未被测量的混杂因子造成的偏倚。这样的例子比比皆是（例如图 6.3），其中向分析中加入一个变量会引入无法弥补的偏倚（Pearl 2009b, 2010a; Shrier 2009; Sjölander 2009）。

激化争议的另一个因素，是一种普遍信念：倾向评分法的偏倚削减潜力可以用实验来评估——即做案例研究，把倾向评分得到的效应估计与受控随机化实验得到的估计相比较（Shadish and Cook 2009）。\footnote{这类信念受到一些豪迈陈述的鼓励，例如：“关于这类分析能够得出与一个完全平行的随机化实验相同结论的戏剧性证据，见 Shadish and Clark (2006, 未发表)”（Rubin 2007）。} 这一信念没有根据，因为倾向评分的偏倚削减潜力关键地取决于 $S$ 的具体选择，或者更准确地说，取决于 $S$ 之内与之外各变量之间的因果关系。在一个问题实例中（比如说俄克拉何马州的一个教育项目）测到显著偏倚，并不排除在另一个实例中（比如说阿肯色州的犯罪控制）测得零偏倚，即便二者的统计分布 $P(x,s,y)$ 完全相同。

鉴于这些考虑，人们有理由提出一个社会科学式的问题：倾向评分究竟有什么特质，妨碍了人们更普遍地理解它的前景与局限？

Richard Berk 在《回归分析：一种建设性批评》（Berk 2004）中回顾了社会科学中类似的现象：完美的想法被科学共同体误读。“我记得与 Don Campbell 的一次谈话，他公开希望自己从未写过 Campbell and Stanley (1966)。那本理应著名的《实验与准实验研究设计》的意图，是把随机化实验与准实验近似相对照，并极力劝阻后者。然而该书的表面影响，却是为一大批准实验设计在广泛的应用社会科学领域中确立了合法性。在我于 Dudley Duncan 职业生涯晚期结识他之后，他说他常常认为他那本有影响力的路径分析著作《结构方程模型导论》是一个大错。研究者读完那本书后相信：关于社会不平等的基本政策问题，可以用路径分析迅速而轻易地得到解答。”（p.~xvii）

我相信，围绕倾向评分也演化出了类似的文化现象。并不是 Rosenbaum 与 Rubin 在陈述成功条件时不够谨慎。形式上，他们非常清楚，告诫实践者倾向评分只在“强可忽略性”条件下才有效。然而他们未能意识到的是：仅仅警告人们提防他们无法识别的危险是不够的；要保护人们免于危险的冒险，我们还必须给他们眼镜去发现威胁，给他们一种有意义的语言去推理这些威胁。由于未能给读者配备工具（例如图）去识别“强可忽略性”如何被违反或如何被达成，他们鼓励了一代研究者（包括联邦机构）去假定：可忽略性要么在多数情形下成立，要么可以通过巧妙的设计使之成立。

\subsection{do 演算背后的直觉}

\qhead{就定理 3.4.1 致作者的问题}

在 do 演算的推理规则中（第 85 页），子图 $G_{\bar X}$ 表示在操作 $\dox{X=x}$ 下通行的分布，因为 $X$ 的一切直接原因都被移除。那么子模型 $G_{\underline X}$ 表示什么分布——在 $X$ 的直接效应被移除之后？

\qhead{作者答复}

图 $G_{\underline X}$ 表示“把 $X$ 的一切子节点保持恒定”这一假设性动作。这切断了从 $X$ 到 $Y$ 的一切有向路径，同时使一切后门路径保持完好。因此，若 $X$ 与 $Y$ 在该图中是 d-连通的，那必定是由于二者之间存在（未被阻断的）混杂路径。反之，如果我们找到一个节点集 $Z$，它在该图中把 $X$ 与 $Y$ d-分离，我们就可以确信 $Z$ 在原图中阻断了所有后门路径。若我们进一步对变量 $Z$ 取条件，那么按后门准则，我们可以确信已中和了一切混杂因子，因而在这样的条件化之后测得的任何相依，都必定来自 $X$ 对 $Y$ 的因果效应，不含混杂。

\subsection{G-估计的有效性}

在 \S 3.6.4 中，我们介绍了 G-估计公式 (3.63) 以及反事实独立条件 (3.62)，
\[
(Y(x)\indep X_k\mid \bar L_k,\bar X_{k-1}=\bar x_{k-1}),
\]
Robins 证明了它是 (3.63) 的充分条件。一般而言，条件 (3.62) 既过强，又缺乏直观基础。导致 (3.63) 的一个更一般、更直观的条件在 (4.5)（第 122 页）中导出，其内容如下：

\begin{quote}
\textbf{（3.62*）g-估计的一般条件（序贯去混杂）}

$P(y\mid g=x)$ 可识别，并由 (3.63) 给出，若从 $X_k$ 到 $Y$ 的每条“避开行动的”后门路径都被 $X_k$ 之非后代的某个子集 $L_k$ 阻断。（所谓“避开行动”，我们指路径中不含任何进入晚于 $X_k$ 的 $X$ 变量的箭头。）
\end{quote}

这一条件相对于 (3.62) 有几处改进，如下面三个例子所示。

\begin{example}\label{ex:11.3.1}
图 11.10 展示了这样的情形：g-公式 (3.63) 在取过去的一个子集 $L_k$ 时有效，而取整个过去时无效。假定 $U_1$ 与 $U_2$ 未被观测，时间顺序为 $U_1,Z,X_1,U_2,Y$，我们看到 (3.62) 与 (3.62*) 在 $L_1=\varnothing$ 时都满足，因而 (3.63) 成立；而取整个过去 $L_1=Z$ 则会违反二者。
\end{example}

\begin{figure}[htbp]
\centering
\begin{tikzpicture}[causal,scale=0.95,every node/.style={transform shape}]
  \node[font=\small] (U1) at (0,1.5) {$U_1$};
  \node[font=\small] (Z) at (1.8,1.5) {$Z$};
  \node[font=\small] (U2) at (3.8,1.5) {$U_2$};
  \node[font=\small] (X1) at (0.15,0.15) {$X_1$};
  \node[font=\small] (Y) at (3.0,0.15) {$Y$};
  \draw[edd] (U1)--(Z);
  \draw[edd] (U2)--(Z);
  \draw[edd] (U1)--(X1);
  \draw[edd] (U2)--(Y);
  \draw[ed] (X1)--(Y);
\end{tikzpicture}
\caption{对整个过去 $L_1=Z$ 取条件会使 g-估计失效。}
\label{fig:11.10}
\end{figure}

\begin{example}\label{ex:11.3.2}
图 11.11 展示了这样的情形：把 $L_k$ 定义为 $X_k$ 的“非后代”集（而非 $X_k$ 的时间前驱）会拓宽 (3.62)。假定时间顺序为 $U_1,X_1,S,Y$，(3.62) 与 (3.62*) 在 $L_1=S$ 时都满足，而在 $L_1=\varnothing$ 时不满足。
\end{example}

\begin{figure}[htbp]
\centering
\begin{tikzpicture}[causal,scale=0.95,every node/.style={transform shape}]
  \node[font=\small] (U1) at (0,1.5) {$U_1$};
  \node[font=\small] (S) at (1.9,1.5) {$S$};
  \node[font=\small] (X1) at (0.15,0.15) {$X_1$};
  \node[font=\small] (Y) at (2.9,0.15) {$Y$};
  \draw[edd] (U1)--(S);
  \draw[edd] (U1)--(X1);
  \draw[ed] (S)--(Y);
  \draw[ed] (X1)--(Y);
\end{tikzpicture}
\caption{把非前驱 $S$ 纳入进来，使 g-估计变得有效。}
\label{fig:11.11}
\end{figure}

\begin{example}\label{ex:11.3.3}
（本例由 Ilya Shpitser 应 Eliezer Yudkowsky 之问而构造。）图 11.12 展示了这样的情形：即便对 $L_k$ 作新的解释，(3.62) 也不满足，而图形条件 (3.62*) 却满足。容易看出 (3.62*) 成立：从 $X_1$ 到 $Y$ 的一切避开行动的后门路径都被 $\{X_0,Z\}$ 阻断。同时可以证明（用孪生网络方法，第 213 页），给定 $Z$ 与 $X_0$ 时，$Y(x_0,x_1)$ 并不独立于 $X_1$。（在孪生网络模型中，存在一条从 $X_1$ 到 $Y(x_0,x_1)$ 的 d-连通路径，如下：
\[
X_1\leftarrow Z\leftrightarrow Z^*\to Y^*.
\]
因此，对 $Y(x_0,x_1)$ 与 $X_1$ 而言，(3.62) 不满足。）
\end{example}

\begin{figure}[htbp]
\centering
\begin{tikzpicture}[causal,scale=1.0,every node/.style={transform shape}]
  \node[font=\small] (X0) at (0,0.9) {$X_0$};
  \node[font=\small] (Z) at (2.0,0.9) {$Z$};
  \node[font=\small] (Y) at (4.2,0.9) {$Y$};
  \node[font=\small] (X1) at (2.9,-0.1) {$X_1$};
  \draw[ed] (X0)--(Z);
  \draw[ed] (Z)--(Y);
  \draw[ed] (Z)--(X1);
  \draw[ed] (X1)--(Y);
  \draw[bid,dashed] (Z) to[bend right=50] (X1);
\end{tikzpicture}
\caption{一个使 g-估计有效、却违反 Robins 条件 (3.62) 的图。}
\label{fig:11.12}
\end{figure}

这个例子再次表明，Robins 最初为导出 (3.63) 所采用的潜在结局语言有其弱点。使 g-估计公式之使用得以正当化的反事实条件 (3.62) 晦涩难懂，唤不起任何直观支持。应用这一公式的流行病学家，是在没有实质医学知识指引的情况下使用它。所幸，图形方法正在迅速进入流行病学实践（Greenland et al. 1999a; Robins 2001; Hernán et al. 2002; Greenland and Brumback 2002; Kaufman et al. 2005; Petersen et al. 2006; VanderWeele and Robins 2007），越来越多的研究者开始理解 g-估计背后的假设。再加上这样一种理解——结构方程模型包摄、统一并支撑着图形、反事实、潜在结局与充分成分（Rothman 1976）诸种因果进路\footnote{这一统一尚未得到领军流行病学家（Greenland and Brumback 2002）、经济学家（Heckman and Vytlacil 2007）与社会科学家的充分强调（Morgan and Winship 2007），更不用说统计学家了（Cox and Wermuth 2004; Rubin 2005）。在对多元主义给予应有尊重的同时，必须指出：一切因果进路都是本书（第七章）所述结构理论的变体或抽象。}——流行病学很有希望成为第一个彻底摆脱昔日教条、“冲破我们在学术上被强化的、不愿直接思考原因的抵触”（Weinberg 2007）的学科。
